\documentclass[10pt]{article}
\usepackage[margin=0.68in]{geometry}
\usepackage[T1]{fontenc}
\usepackage[utf8]{inputenc}
\usepackage{booktabs}
\usepackage{array}
\usepackage{tabularx}
\usepackage[hidelinks]{hyperref}
\setlength{\parindent}{0pt}
\setlength{\parskip}{0.35em}
\setlength{\emergencystretch}{2em}

\newcommand{\questionsep}{\par\vspace{1.1em}}

\begin{document}

\begin{center}
{\Large Response to Reviewer kUmZ}\\[2pt]
Submission 13175
\end{center}

\textbf{Opening response.}
We thank the reviewer for identifying the central issues: molecular-target
provenance, experimental breadth, and what cell alignment adds beyond spot
prediction. We now distinguish native cell-resolved measurements from derived
Visium HD targets, report all 25 release categories, and retain weak transfer
results alongside positive findings.

\questionsep
\textbf{Q1 (P2, Q2, Q4, L1--L2).}\quad
\textit{Are the molecular targets truly single-cell?}

\textbf{A1 (Molecular-target provenance).}\enspace
The resource contains three evidence types:

\begin{center}
\begingroup
\footnotesize
\setlength{\tabcolsep}{3.5pt}
\renewcommand{\arraystretch}{1.08}
\begin{tabularx}{\linewidth}{
  >{\raggedright\arraybackslash}p{0.18\linewidth}
  >{\raggedright\arraybackslash}X
  >{\raggedright\arraybackslash}p{0.27\linewidth}
  >{\raggedright\arraybackslash}p{0.20\linewidth}}
\toprule
Component & Evidence unit & Count & Status \\
\midrule
Public Xenium & Platform-segmented cells and transcripts &
16,314,129 cells & Native cell-resolved \\
Public Visium HD & Native 2-\(\mu\)m bins aggregated to cell masks &
7,629,697 derived cells &
Derived cell-aligned \\
New in-house DBiC-seq & Paired cell morphology and RNA &
53,989 post-QC cells from 21 samples & Native cell-resolved \\
\bottomrule
\end{tabularx}
\endgroup
\end{center}

\begin{itemize}
\item All non-HD data are native cell-resolved measurements: Xenium and
DBiC-seq assign an expression profile to each segmented cell. Only Visium HD
uses bin-to-cell aggregation, following 10x Genomics'
\href{https://www.10xgenomics.com/support/software/space-ranger/latest/analysis/segmented-outputs}{officially
documented mapping} from native 2-\(\mu\)m barcodes to segmented cells. The new
in-house DBiC-seq data are available for
\href{https://anonymous.4open.science/r/sMMC-22M-DB75/README.md}{download and
preview}.
\item We will therefore not describe the entire collection as directly measured
single-cell data. We will consistently use ``native cell-resolved'' for
Xenium/DBiC-seq and ``derived cell-aligned'' for Visium HD, and retitle the work
\emph{sMMC: A Cell-Aligned Multimodal Resource for Spatial Transcriptomics}.
\end{itemize}

\textbf{A2 (Native-cell validation).}\enspace
To directly address this concern, we reran the benchmark using only native
Xenium data; Table~1 reports the 25-organ results. Table~2 separately reports
cross-sample organ transfer on Visium HD.

\newpage
\questionsep
\textbf{Q2 (P3, Q1).}\quad
\textit{Are all 25 release categories evaluated?}

\textbf{A1 (25-organ results).}\enspace
Table~1 reports results for all 25 organs under one fixed protocol.

\begin{center}
\begingroup
\footnotesize
\setlength{\tabcolsep}{4.2pt}
\renewcommand{\arraystretch}{1.04}
\textit{\textbf{Table 1. Per-organ benchmark results.}
Results use the top-50 training-selected HVGs, four contiguous spatial
holdouts, and a 5\% train--test buffer.}\\[5pt]
\begin{tabular*}{\linewidth}{@{\extracolsep{\fill}}lrrrrrr@{}}
\toprule
& \multicolumn{3}{c}{Gene Pearson} & & & \\
\cmidrule(lr){2-4}
\shortstack[l]{Release category\\(evaluated species)} &
Image &
\shortstack{Coordinate\\only} &
\shortstack{Spatial\\KNN} &
\shortstack{Gene\\Spearman} &
\shortstack{Cell\\Pearson} &
F1 \\
\midrule
\textbf{Bone (mouse)} & 0.030 & \textbf{0.124} & 0.055 & 0.037 & 0.180 & 0.240 \\
Brain (human + mouse) & \textbf{0.399} & 0.010 & 0.036 & 0.371 & 0.570 & 0.752 \\
Breast (human) & \textbf{0.400} & 0.053 & 0.035 & 0.375 & 0.433 & 0.487 \\
Cervical (human) & \textbf{0.178} & 0.017 & 0.014 & 0.157 & 0.247 & 0.026 \\
\textbf{Colon (mouse)} & \textbf{0.615} & -0.137 & 0.072 & 0.588 & 0.739 & 0.744 \\
Colorectal (human) & \textbf{0.431} & 0.056 & 0.083 & 0.419 & 0.532 & 0.623 \\
\textbf{Embryo (mouse)} & \textbf{0.278} & 0.037 & 0.045 & 0.251 & 0.517 & 0.342 \\
Head (zebrafish) & \textbf{0.216} & 0.022 & 0.031 & 0.193 & 0.429 & 0.287 \\
Heart (human) & \textbf{0.292} & 0.010 & 0.021 & 0.266 & 0.688 & 0.304 \\
Kidney (human) & \textbf{0.402} & 0.019 & 0.017 & 0.361 & 0.471 & 0.460 \\
Liver (human) & -0.002 & 0.003 & \textbf{0.010} & -0.002 & 0.561 & 0.400 \\
Lung (human) & \textbf{0.359} & 0.065 & 0.053 & 0.317 & 0.402 & 0.457 \\
Lymph Node (human) & \textbf{0.141} & 0.043 & 0.013 & 0.131 & 0.164 & 0.023 \\
Ovarian (human) & \textbf{0.250} & 0.122 & 0.104 & 0.239 & 0.312 & 0.234 \\
Ovarian glands (human) & \textbf{0.402} & 0.059 & 0.043 & 0.391 & 0.559 & 0.756 \\
Pancreas (human) & \textbf{0.324} & 0.090 & 0.068 & 0.272 & 0.505 & 0.275 \\
Pancreatic (human) & \textbf{0.366} & 0.080 & 0.051 & 0.341 & 0.540 & 0.694 \\
Pancreatic duct gland (human) & \textbf{0.331} & 0.091 & 0.100 & 0.282 & 0.408 & 0.386 \\
Plant (\textit{A. thaliana}) & -0.011 & 0.004 & \textbf{0.012} & -0.009 & 0.385 & 0.052 \\
Prostate (human) & \textbf{0.263} & -0.015 & 0.025 & 0.244 & 0.263 & 0.083 \\
Seed (soybean) & -0.020 & -0.003 & \textbf{0.008} & -0.018 & 0.314 & 0.037 \\
Skin (human) & \textbf{0.359} & 0.048 & 0.086 & 0.299 & 0.437 & 0.368 \\
\textbf{Small Intestine (mouse)} & \textbf{0.572} & -0.104 & 0.067 & 0.548 & 0.701 & 0.715 \\
Tonsil (human) & \textbf{0.294} & 0.028 & 0.018 & 0.249 & 0.462 & 0.399 \\
Xenograft (human + mouse) & \textbf{0.387} & 0.041 & 0.049 & 0.362 & 0.526 & 0.589 \\
\midrule
\textbf{30-sample macro} & \textbf{0.324} & 0.046 & 0.053 &
\textbf{0.291} & \textbf{0.442} & \textbf{0.413} \\
\bottomrule
\end{tabular*}
\endgroup
\end{center}

\textbf{A2 (Cross-sample transfer).}\enspace
We additionally evaluated sample-to-sample transfer across eight organs.
\begin{center}
\begingroup
\footnotesize
\setlength{\tabcolsep}{3.4pt}
\renewcommand{\arraystretch}{1.06}
\textit{\textbf{Table 2. Eight-organ cross-sample Visium HD benchmark
(source$\rightarrow$target sample).} TM is training mean; its Gene Pearson is
undefined.}\\[3pt]
\begin{tabular*}{\linewidth}{@{\extracolsep{\fill}}llrrrrr@{}}
\toprule
Organ & Pair & \shortstack{UNI2-h\\Gene P} & \shortstack{UNI2-h\\Cell P} &
\shortstack{UNI2-h\\F1} & \shortstack{TM\\Cell P} & \shortstack{TM\\F1} \\
\midrule
Human breast & j$\rightarrow$m & 0.0272 & 0.1109 & 0.1794 & 0.1108 & 0.1215 \\
Human ovary & ad$\rightarrow$ak & 0.0267 & 0.4593 & 0.1161 & 0.4906 & 0.0561 \\
Human lung & k$\rightarrow$d & 0.0143 & 0.2121 & 0.0960 & 0.1455 & 0.0236 \\
Human pancreas & g$\rightarrow$ag & 0.0040 & 0.0090 & 0.0389 & 0.0023 & 0.0168 \\
Human tonsil & u$\rightarrow$i & 0.0141 & 0.2769 & 0.0708 & 0.3264 & 0.0346 \\
Mouse brain & e$\rightarrow$ah & 0.0243 & 0.2016 & 0.0623 & 0.2238 & 0.0179 \\
Mouse embryo & b$\rightarrow$ai & 0.0048 & 0.0913 & 0.0369 & 0.1683 & 0.0134 \\
Mouse kidney & a$\rightarrow$aj & 0.0049 & 0.2678 & 0.0513 & 0.4696 & 0.0159 \\
\midrule
\textbf{Organ macro} & --- & \textbf{0.0151} & \textbf{0.2036} &
\textbf{0.0815} & \textbf{0.2422} & \textbf{0.0375} \\
\bottomrule
\end{tabular*}
\endgroup
\end{center}

\begin{itemize}
\item The cross-sample Gene Pearson is indeed low (organ macro 0.0151), and
UNI2-h improves F1 but not Cell Pearson over the training mean. We do not claim
to have solved this setting.
\item The low Pearson is not evidence that the dataset is invalid: the in-domain
benchmark and separate construction audit establish usable targets. Instead,
this negative result exposes the still-unsolved difficulty of single-cell
generalization under patient, platform, composition, acquisition, and
morphology shifts.
\item The dataset makes this failure measurable and provides a basis for future
solutions. We will use ``cross-patient'' only for verified donors and otherwise
say ``sample-held-out.''
\end{itemize}

\questionsep
\textbf{Q3 (Q4).}\quad
\textit{What techniques were used to interpolate or match spot-level gene
expression to individual cells?}

\textbf{A.}
\begin{itemize}
\item This premise does not apply to our main benchmark: its input is histology
alone and its target is a cell-resolved RNA profile. No spot-level
gene-expression measurement is supplied to the model. Training directly pairs
each cell's RNA profile with two cell-centered histology crops: a local crop
around the cell and a larger crop capturing tissue context.
\item STBoost is therefore not a spot-to-cell expression-interpolation method.
It adapts existing image-to-spot predictors to cell-level prediction by
replacing the spot-centered image input with these two hierarchical
cell-centered crops, fusing their representations, and retaining the
corresponding spot-level method's downstream prediction modules.
\item Because this is primarily a dataset paper, we placed the full architecture
and training details in the Appendix. We agree that the distinction was not
sufficiently clear and will add a concise description of the task input,
target, training, and inference procedure to the main body.
\end{itemize}

\questionsep
\textbf{Q4 (P5, P6, minor remarks).}\quad
\textit{Was Figure 1 AI-generated, and how will the writing and
figure-presentation issues be addressed?}

\textbf{A1 (Figure 1 provenance).}\enspace
Figure~1 was not AI-generated. I began collecting this resource in the first
year of my PhD and have continued developing it for several years. I manually
drew and revised two versions of Figure~1 in Sketch; together they contain 162
editable layers and required more than one week of work. Screenshots documenting
the editable layers are available in our
\href{https://anonymous.4open.science/r/sMMC-22M-DB75}{anonymous GitHub
repository} under \texttt{figure1\_design\_evidence/}.

\textbf{A2 (Presentation revisions).}\enspace
We thank the reviewer for identifying these issues and will address them
individually:
\begin{itemize}
\item correct the capitalization, missing letter/space, and blank elements in
Figure~1, and expand its caption;
\item define ``study split,'' ``cell unit,'' and the Figure~2b training ratio at
first use;
\item remove the Figure~3 border, render its labels as vector text, make the
requested line-level cuts, and state precisely which parts of the resource are
publicly released.
\end{itemize}

\questionsep
\textbf{Q5 (P4, P9).}\quad
\textit{What can the context analysis support, and how complete is the
demographic metadata?}

\textbf{A1 (Context claim).}
\begin{itemize}
\item The ovarian AK/AD/AL comparison cannot establish an age effect because
age is nested within sample/patient and panels differ. We will rename it
``sample/age-confounded context shift.''
\item Its role is to motivate
Average/Worst/Gap/Support auditing, not causal attribution.
\end{itemize}

\textbf{A2 (Metadata boundary).}
\begin{itemize}
\item Donor, age, sex, and disease are incompletely reported across sources, and
ethnicity is undocumented.
\item We will report missingness, never infer sensitive attributes, and analyze
subgroups only with adequate independent-donor support.
\end{itemize}

\questionsep
\textbf{Q6 (P5, P7, minor comments).}\quad
\textit{Please clarify terminology and correct the remaining presentation issues.}

\textbf{A1 (Terminology).}
\begin{itemize}
\item STBoost is the model-agnostic cell-aligned interface.
\item STBoosted BLEEP is published BLEEP retrained through it.
\item STBoost-Ref is our image-only retrieval reference predictor; at inference
it receives local/context histology, never query expression.
\end{itemize}

\textbf{A2 (Presentation and limitations).}
\begin{itemize}
\item We will use these names consistently, define every split/unit, and correct
Figure~1 and Figure~3.
\item We will state the limits of spatial autocorrelation, sample/platform shift,
HD boundary assignment, metadata missingness, and prediction as a substitute
for measurement.
\end{itemize}

\questionsep
\textbf{Q7 (Q3, L3).}\quad
\textit{What does cell-aligned evaluation add beyond spot-centered prediction?}

\textbf{A1 (Matched resolution control).}
\begin{itemize}
\item We constructed 8-, 16-, and 55-\(\mu\)m pseudo-spots from the same
native-Xenium cells while holding the image representation, target genes, and
evaluation protocol fixed across six samples.
\item Mean Gene Pearson was 0.365 at cell level, 0.365 at 8 \(\mu\)m, 0.363 at
16 \(\mu\)m, and 0.330 at 55 \(\mu\)m. We therefore do not claim that
cell-level prediction is universally more accurate than spot-level prediction.
\end{itemize}

\textbf{A2 (What averaging hides).}
\begin{itemize}
\item In full-density lung HLCX022, 55-\(\mu\)m aggregation mixed cell types in
55.5--66.4\% of test pseudo-spots, containing 73.8--81.0\% of evaluated cells.
In the much sparser HHDX011 sample, only 3.5--4.1\% of pseudo-spots were mixed,
containing 7.0--8.3\% of cells.
\item The benefit of cell alignment is therefore density dependent: it
preserves within-region cellular heterogeneity that coarse averaging can
conceal, even when aggregate correlation changes little. We will replace the
vague limitation identified by the reviewer with this measured statement.
\end{itemize}

\questionsep
\textbf{Q8 (P6).}\quad
\textit{How are cells detected and aligned when the platform does not provide
cell boundaries, and how reliable is the resulting Visium HD assignment?}

\textbf{A1 (Platform-specific construction).}
\begin{itemize}
\item Xenium supplies native cell boundaries and transcript coordinates; no
spot-to-cell expression inference is used.
\item For Visium HD, official spatial transforms register native
2-\(\mu\)m bins to H\&E. A published CellViT model supplies same-frame cell
contours; intersecting bins are aggregated within each contour, conflicts are
assigned to the nearest full-resolution cell centroid, and unsupported polygons
are excluded. Thus there is no single spot center that is simply matched to a
cell center.
\end{itemize}

\textbf{A2 (Assignment sensitivity).}\enspace
We audited 3,000 raw CellViT polygons in each of three Visium HD samples
(9,000 total) under \(\pm 1\)-\(\mu\)m registration shifts:
\begin{center}
\begingroup
\footnotesize
\setlength{\tabcolsep}{3.2pt}
\renewcommand{\arraystretch}{1.07}
\begin{tabular*}{\linewidth}{@{\extracolsep{\fill}}lrrrr@{}}
\toprule
Visium HD sample & \shortstack{Polygons with\\canonical bins} &
\shortstack{Shift-bin\\Jaccard} &
\shortstack{Shift-expression\\cosine} &
\shortstack{Shift median absolute\\UMI change} \\
\midrule
Human lung cancer & 49.4\% & 0.727--0.733 & 0.954--0.957 & 11.1--11.3\% \\
Mouse brain & 97.5\% & 0.806--0.816 & 0.936--0.939 & 6.0--6.1\% \\
Human pancreas & 50.0\% & 0.706--0.714 & 0.994 & 13.0--13.7\% \\
\bottomrule
\end{tabular*}
\endgroup
\end{center}
\begin{itemize}
\item Among supported polygons, small shifts largely preserve expression
direction, but coverage is sample dependent. Mask erosion/dilation is more
disruptive (Jaccard 0.462--0.720; expression cosine 0.871--0.993; median
absolute UMI change 32.8--66.6\%).
\item We will move the segmentation, transform, conflict-resolution, filtering,
and sensitivity details into the main body; release the corresponding
coverage/QC fields; and continue to report native Xenium separately from
derived HD evidence.
\end{itemize}

\newpage
\questionsep
\textbf{Q9 (P8, P9, L4).}\quad
\textit{How many independent patients or animals support the experiments, and
how complete is the demographic metadata?}

\textbf{A1 (Verified evaluation units).}
\begin{center}
\begingroup
\footnotesize
\setlength{\tabcolsep}{4pt}
\renewcommand{\arraystretch}{1.07}
\begin{tabularx}{\linewidth}{@{}p{0.27\linewidth}p{0.25\linewidth}X@{}}
\toprule
Evaluation & Verifiable sample support & Patient/animal interpretation \\
\midrule
Native-Xenium breadth benchmark & 30 H\&E-aligned samples &
Independent donor/animal IDs are not consistently reported upstream \\
Eight-organ HD transfer & 16 source/target samples in 8 pairs &
Sample-held-out; donor relationships are unverified \\
\bottomrule
\end{tabularx}
\endgroup
\end{center}
\begin{itemize}
\item We cannot responsibly convert sample records into a unique patient or
animal count when upstream identifiers are absent. We will add per-organ
sample-record counts alongside cell counts, mark unavailable donor/animal
fields explicitly, and use ``patient-held-out'' only where independent donor
identity is verified.
\end{itemize}

\textbf{A2 (Demographic coverage).}
\begin{itemize}
\item As stated in Q5, age, sex, and disease are retained only when reported by
the source, while ethnicity is undocumented in the current manifest. We will
report field-level missingness rather than infer sensitive attributes, and no
demographic comparison will be presented as independently supported without
multiple verified donors per group.
\end{itemize}

\questionsep
\textbf{Q10 (P7, second Q4).}\quad
\textit{How will STBoost and the changing resource size be named?}

\par\textbf{A.}
\begin{itemize}
\item We will define STBoost at its first appearance in both the abstract and
Introduction, introduce BLEEP before the first comparison, and use STBoost,
STBoosted BLEEP, and STBoost-Ref consistently instead of the undefined label
``Ours.''
\item To avoid 22M/23M becoming obsolete as the resource grows, we will use the
count-independent name \textit{sMMC} and report the exact frozen release version
and platform-resolved counts in the manifest.
\end{itemize}

\questionsep
\textbf{Q11 (P3--P5, Q1).}\quad
\textit{How will the revised main body distinguish the three experimental
purposes and reduce redundancy?}

\par\textbf{A.}
\begin{itemize}
\item The resolution evidence---native-Xenium targets, the 25-category
benchmark, and the matched pseudo-spot control---addresses whether cell-aligned
evaluation preserves cell-specific molecular heterogeneity.
\item The context experiment is orthogonal: it tests whether a trained model
remains reliable across recorded biological contexts; it is not offered as
evidence that the targets are single-cell.
\item We will place a compact 25-category result summary, the histology-only
input and cell-expression target, split definitions, and platform-specific
alignment/QC in the main body, while removing repeated
scale/resolution/context framing and the line-level redundancies identified by
the reviewer.
\end{itemize}

\end{document}
